\section{Memory limits and larger collections}
\subsection{Put the stages back together}
We can now describe one complete request. Let $T_{\rm shared}$ include query embedding,
normalization, routing, and any common filtering. Let $T_a^{\rm local}$ be the local cost for
A, B, or C derived earlier, including its score table. Then
\begin{equation}
T_a^{\rm request}=T_{\rm shared}+T_a^{\rm local}+T_{\rm final},
\label{eq:total}
\end{equation}
where $T_{\rm final}$ contains output work and any enabled reranking. This equation sums stages
for the same request. It does not justify adding unrelated stage medians or adding the real
encoder's measured time to constructed vectors that the encoder never produced.

The corresponding memory summary separates stored allocation from live query scratch.
Write $M_0$ for allocated document rows, routing state, and common containers. Let $D_C$ be C's
allocated directory, $Q_B$ its counterpart B's extra masks and queue, and $Q_C$ C's extra key
queue. The common query scratch remains $Q_0$.
\begin{center}
\begin{tabular}{@{}lll@{}}\toprule
Method & Stored index allocation & Query scratch\\\midrule
A & $M_0$ & $Q_0$\\
B & $M_0+M_{\rm planes,allocated}$ & $Q_0+Q_B$\\
C & $M_0+D_C$ & $Q_0+Q_C$\\\bottomrule
\end{tabular}
\end{center}
Actual common allocation can differ with container capacities. Add the runtime and other
resident data once when applying a process RAM limit. Do not add overlapping process and GPU
accounting blindly. A build-memory limit must also include temporary input and output copies.

\subsection{Start with the memory calculation at one billion rows}
At $N=10^9$, $d=256$, and eight-byte IDs, packed codes occupy 32 GB and IDs occupy 8 GB.
These two arrays alone exceed a 32 GB budget. B keeps a second sign representation, adding
at least 32 GB before per-cluster padding. Changing the number of probes cannot remove this
stored payload.

\begin{center}
\begin{tabular}{@{}lrrl@{}}\toprule
Layout at 1B rows & Payload floor & Planned query RAM & 64 GB budget\\\midrule
A, direct prefix route & 40 GB & about 49 GB & Modeled fit\\
C, global 13-bit key & 40 GB & about 49 GB & Modeled fit\\
B, global cluster & 72 GB & about 83 GB & Rejected\\
B, prefix route, large leaf & 72 GB & about 81 GB & Rejected\\\bottomrule
\end{tabular}
\end{center}
The planning values include capacity and container allowances and a configurable 1 GB runtime
reserve; the reserve is not measured runtime memory. All three fail 32 GB in their current
in-memory layouts. All listed planning cases fit 1 TB, but a machine with 1 TB need not have
this machine's memory performance. Compressed IDs or another storage architecture would require
a new calculation, not a different label on the same index.

\subsection{When a fixed prefix carries the useful information}
Under the earlier independent strong-coordinate model, thirteen fixed strong coordinates give
an expected ideal group of $10^9/2^{13}\approx122{,}070$ rows. This is far more than the requested
100. The distributional calculation controls the small probability of an insufficient group;
it is not a claim about ordinary text embeddings.

A can route directly to this group and scan it. C can look up its key and score the same rows.
Both cost roughly 3 ms in the calibrated conditional model. B can also use this route and a
large leaf, producing a similar time with more stored data. We should not force B through a
global billion-row bitmap just to make C appear better. For this case, A and C are closely
related ways to obtain the same candidate set, and the model cannot resolve a few-percent gap.

\subsection{When the important coordinates change with each query}
A fixed prefix may now contain little useful signal. B can choose its split coordinates from
the query. Keeping thirteen useful cuts and scaling the leaf size to 160,000 preserves the
constructed path. Its last-split mask payload is approximately
\[
(13+2)\ceil{10^9/64}\times8=1.875\ \text{GB}.
\]
The factor 15 counts the simultaneously live path and alternative masks for this prescribed
path, including parent/child overlap. Other search histories can need different scratch.

The complete native size model predicts about 0.10 s for this global B path and about 25 s for
a global scan. Doubling B's variable cost gives about 0.20 s; this is a sensitivity example,
not a confidence interval. No full billion-document query was measured.

A useful comparison gives A the benefit of routing. If it scores fraction $f$ of the collection,
its approximate cost is $T_{\rm route}+fNc_{\rm scan}$. Even ignoring routing time, A must
score less than about 0.41\% of the corpus to undercut the 0.10 s B prediction. It must also
meet the recall target at that fraction. If routing produces a small, accurate candidate set
cheaply, scanning can win. If it cannot, the query-dependent B path is worth testing.

For C, a short key must actually capture important coordinates. In the constructed changing
case, a fixed 13-coordinate key misses all thirteen strong coordinates with probability about
49.94\%. A nonempty directory lookup can therefore still leave almost the entire scoring task.
Key occupancy alone cannot establish the latency or recall advantage.

\subsection{What we can conclude}
On the measured real embeddings, the scan remains the supported practical baseline. Under a
strong fixed-prefix condition, direct routing and key lookup both avoid most scores. Under
changing strong coordinates, B's query-dependent splits produced a large measured advantage
in constructed data. Each conclusion has a stated workload and a tested parameter space.

The formulas help explain those outcomes and suggest settings. They do not prove that every
possible IVF layout loses, that Matryoshka creates the constructed distribution, or that a
local timing coefficient remains exact at a billion rows. The next useful experiment would
check whether a real representation has the score concentration and routing behavior required
by the favorable case, while preserving the same recall and RAM constraints.
